\section{Conclusion}
\label{sec:conclusion}

All major published RLEF papers for code LLMs train with binary pass/fail reward \citep{le2022coderl,shojaee2023ppocoder,gehring2024rlef,yu2025dapo}. We showed that this choice produces zero gradient contribution in 98.7\% of early-training GRPO groups under realistic partial-pass distributions --- a dead zone that ratio reward ($k/n$ test cases passing) mathematically escapes. The mechanistic guarantee is exact: whenever a GRPO group contains at least one completion with partial credit and no completion with full credit, binary reward must produce zero advantage variance while ratio reward must produce non-zero advantage variance (0.0475 vs.\ 0.0000 for the illustrative group; 987/1,000 simulation groups).

We confirmed this guarantee experimentally (h-e1) and tested its policy-level implications (h-m1). The h-m1 experiment produced a precision null result: both conditions showed \texttt{reward\_mean = 0.0} at all 208 logged training steps, with \texttt{clipped\_ratio = 1.0} throughout --- precisely confirming that \texttt{max\_new\_tokens = 512} is insufficient for 6.7B models to produce executable solutions for APPS competition-level problems, making both reward functions mathematically equivalent throughout training.

The next step is clear. Re-run h-m1 with \texttt{max\_new\_tokens = 1,024--2,048} on APPS, or switch training data to HumanEval or MBPP where DeepSeek-Coder-6.7B achieves 40--65\% base pass@1. Under those conditions, the mechanistic guarantee we proved will be active, and the policy-level prediction becomes testable. The validated infrastructure (35/35 unit tests, 9 engineering issues resolved, FractionPartialCallback diagnostic) is ready.

Reward granularity in RLEF is not a minor implementation detail. For models trained with GRPO on hard datasets, the choice between binary and ratio reward determines whether the model receives any learning signal during the critical early phase of training. We have provided the mathematical foundation for this claim. The empirical validation, under correctly-powered conditions, is the natural next step.

\textbf{Future Work.} Redesign h-m1 with longer generation or easier training data to test policy-level effects under the correct prerequisites. Once h-m1 is validated, test the downstream hypotheses (training-evaluation reward alignment h-m2; cross-benchmark generalization h-m3). Long-term: derive a reward granularity selection criterion based on model base solve rate at training time.
